\subsection{Shared Extended Descriptor}
A compact extension form selects a descriptor family and its exact byte length. The descriptor bytes select an EA mode
within that family. For EXT1 and EXT2, the compact form also selects the displacement width.
Only compact forms select extensions; an extended descriptor cannot select another extension.
Field roles are unique across the compact form and its descriptor; payload roles are likewise unique across their selected encoding variants.
Auto-updates declared by either selected encoding are retained. If both encodings declare an auto-update,
their distinct base and index roles are evaluated in base-before-index order, regardless of which encoding declares each role.

An extended memory mode selects its segment from the encoded SREG (:term:base.terms.field:) SEG(s), the fixed SS segment for SP, the fixed CS
segment for PC, or the operation\textquotesingle{}s default data segment.

For a compact displacement, absolute address, or immediate, the selected (:term:base.terms.payload:) follows the compact (:term:base.terms.field:) at that
operand\textquotesingle{}s (:term:base.terms.payload:) position. Extended-descriptor construction begins with the compact escape value in the (:term:base.terms.opcode:)
space and the descriptor bytes declared by that compact form. EXT1 declares one byte and EXT2 declares two bytes.
A displaced form places its selected displacement after
all extended descriptor bytes in the instruction. When an instruction has more than one extended EA operand, the decoder consumes
their complete descriptors in declared instruction operand order. It then consumes every remaining appended EA (:term:base.terms.payload:),
including any selected displacement, in declared instruction operand order.

\clearpage
